% ---
% id: statistics/binomial-box/prop:bracketed-endpoints
% kind: prop
% title: "The enclosure evaluation of the exact rule"
% summary: >-
%   Endpoints bracketed by certified tail enclosures form an admissible rule
%   whose box contains the exact rule's. Each equals the exact quantile where
%   its bracket is decided; otherwise, where the tail decays at the quantile,
%   the excess in counts is bounded through the enclosure's slack and that
%   decay.
% prominence: standard
% status: implemented
% uses: [def:endpoint-rule, def:interval, prop:rule-exact, lem:guard, lem:tail-majorant, lem:pmf-enclosure, lem:log-enclosure, def:construction, statistics/def:conservative-evaluation]
% mentions: [rem:local-evaluation, thm:monotonicity-limits, req:rule-in-force, req:search-procedure]
% assumes: []
% refs: []
% provenance: newly-derived
% implements: certiqs_key.statistics.enclosure:bracketed_upper_endpoint
% lean:
% notes: "Monotonicity in the allocation is not claimed for this evaluation. The edge to the theorem on monotonicity and limiting cases is a mention, not a uses edge: the lead-in points at eq:width, and no clause of this proposition rests on that theorem."
% ---

The \tref[exact binomial quantile rule]{exact rule} compares binomial tails
with their budgets. Where exact sums are too costly, the comparisons can be
decided against enclosures of the tails instead, which is what a
\tref{conservative evaluation} (Definition~\ref{def:conservative-evaluation})
admits: a comparison decided against an upper bound on the tail. A request selects this evaluation
(Requirement~\ref{req:rule-in-force}), and its endpoints are found by the one
search procedure Requirement~\ref{req:search-procedure} fixes;
Remark~\ref{rem:local-evaluation} discusses the search, the enclosures and
their cost. Through the closing sentences of (iii), the width bound
\eqref{eq:width} of Theorem~\ref{thm:monotonicity-limits}\,(iv), proved for
the exact rule, carries over to this evaluation wherever the bracket is
decided and, where it is not, up to the excess (iii) bounds, divided by $n$;
beyond that it is not asserted for this evaluation.

\begin{proposition}\label{prop:bracketed-endpoints}
Assume a \tref{conservative evaluation}
(Definition~\ref{def:conservative-evaluation}). Let $n\ge1$, let
$p\in(0,1)$ and $\varepsilon^{+}\in(0,1)$ be rational, let
$\Ncount{}$ be binomial with parameters $(n,p)$, and let $\kappa^{+}$ be the
exact upper quantile of Definition~\ref{def:interval} at
$(n,p,\varepsilon^{+})$. Let
\begin{equation}\label{eq:tail-enclosure}
  \tau^{-}(k)\ \le\ \Pr[\Ncount{}>k]\ \le\ \tau^{+}(k),\qquad k=0,\dots,n,
\end{equation}
be an enclosure with $\tau^{-}(n)=\tau^{+}(n)=0$, each bound either a
rational or given through a rational bound on its logarithm, such as the
one Lemma~\ref{lem:tail-majorant}\,(ii)--(iii) yields with the mass point
enclosed by Lemma~\ref{lem:pmf-enclosure} and the logarithms by
Lemma~\ref{lem:log-enclosure}. Call $k$ \emph{certified admissible},
$\mathsf{C}^{+}(k)$, when the available upper bound proves
$\Pr[\Ncount{}>k]\le\varepsilon^{+}$: $\tau^{+}(k)\le\varepsilon^{+}$, or in
log space $U(k)\le\lambda^{-}$ for rationals $U(k)\ge\ln\Pr[\Ncount{}>k]$
and $\lambda^{-}\le\ln\varepsilon^{+}$; and \emph{certified inadmissible},
$\mathsf{D}^{+}(k)$, when the lower bound proves
$\Pr[\Ncount{}>k]>\varepsilon^{+}$: $\tau^{-}(k)>\varepsilon^{+}$, or in log
space $\Lambda(k)>\lambda^{+}$ for rationals
$\Lambda(k)\le\ln\Pr[\Ncount{}>k]$ and $\lambda^{+}\ge\ln\varepsilon^{+}$.
An index that is neither is \emph{undecided}. The predicate
$\mathsf{C}^{+}(n)$ always holds.

A \emph{bracketed upper endpoint} is an integer $\tilde k^{+}\in\{0,\dots,n\}$
with $\mathsf{C}^{+}(\tilde k^{+})$ and, if $\tilde k^{+}\ge1$,
$\lnot\mathsf{C}^{+}(\tilde k^{+}-1)$. The \emph{enclosure evaluation of the
exact rule} is the endpoint rule
$R_{\mathrm{encl}}$ whose $R^{+}_{\mathrm{encl}}$ is the bracketed upper
endpoint returned by a fixed deterministic search procedure that returns such
an endpoint for every input, whose $R^{-}_{\mathrm{encl}}$
is obtained in the same way for the lower tail through the mirror
$\Ncount{}'=n-\Ncount{}$ of Lemma~\ref{lem:tail-majorant}\,(v), and which
returns the values of Definition~\ref{def:endpoint-rule}\,(b) and (c) at a
zero budget and at $p\in\{0,1\}$. Then:
\begin{enumerate}[label=(\roman*),leftmargin=2.2em]
\item \emph{(admissible relaxation)} Every bracketed upper endpoint satisfies
$\Pr[\Ncount{}>\tilde k^{+}]\le\varepsilon^{+}$ and
$\kappa^{+}\le\tilde k^{+}$; mirror-wise every bracketed lower endpoint
$\tilde k^{-}$ satisfies $\Pr[\Ncount{}<\tilde k^{-}]\le\varepsilon^{-}$ and
$\tilde k^{-}\le\kappa^{-}$. Hence $R_{\mathrm{encl}}$ is an
\tref{admissible endpoint rule}, Lemma~\ref{lem:guard} applies to it, and the
box
$\mathsf{Acc}_{R_{\mathrm{encl}}}$ emits (Definition~\ref{def:construction})
contains the box the exact rule emits: every $q^{-}$ is at most and every
$q^{+}$ at least the exact rule's.
\item \emph{(the bracket decided)} If $\tilde k^{+}\ge1$ and
$\mathsf{D}^{+}(\tilde k^{+}-1)$ holds, then $\tilde k^{+}=\kappa^{+}$; if
$\tilde k^{+}=0$ then $\tilde k^{+}=\kappa^{+}$. The mirror statement holds
for the lower side.
\item \emph{(conservatism in counts)} Suppose $\tilde k^{+}>\kappa^{+}$,
$\kappa^{+}\le n-2$ and $r_{\kappa^{+}+1}<1$, with $r_j$ the ratios
\eqref{eq:pmf-ratio}. Let $w\ge0$
be the slack of the failed certificate at $\tilde k^{+}-1$:
$w\coloneqq\ln\tau^{+}(\tilde k^{+}-1)-\ln\Pr[\Ncount{}>\tilde k^{+}-1]$, or
in log space
$w\coloneqq U(\tilde k^{+}-1)-\ln\Pr[\Ncount{}>\tilde k^{+}-1]
+(\ln\varepsilon^{+}-\lambda^{-})$. Then
\begin{equation}\label{eq:bracket-conservatism}
  \tilde k^{+}-\kappa^{+}\ <\ 1+\frac{w}{-\ln r_{\kappa^{+}+1}} .
\end{equation}
The emitted endpoint is therefore the exact quantile up to a number of
counts bounded by the enclosure's slack divided by the per-count decay of
the tail at the quantile. The mirror statement holds for the lower side,
through $\Ncount{}'$ (Lemma~\ref{lem:tail-majorant}\,(v)): with
$\kappa^{-}=n-\kappa'^{+}$ and $\tilde k^{-}=n-\tilde k'^{+}$, where the
primed quantities are those of $\Ncount{}'$ at level $\varepsilon^{-}$,
$\kappa^{-}-\tilde k^{-}<1+w'/(-\ln r'_{\kappa'^{+}+1})$ under the mirrored
hypotheses. Consequently, wherever the bracket is decided the emitted
endpoint is the exact rule's, and so are its guarded endpoint and its box
bound; wherever it is undecided the guarded endpoint exceeds the exact rule's
by at most $\tilde k^{+}-\kappa^{+}$ counts and the box bound $q^{+}$ the
exact rule's by at most $(\tilde k^{+}-\kappa^{+})/n$, which
\eqref{eq:bracket-conservatism} bounds under the hypotheses above;
mirror-wise for $k^{-}$ and $q^{-}$. Where the hypothesis $r_{\kappa^{+}+1}<1$ fails (the
exact quantile at or below the mean, reachable at a near-unit tail budget,
Proposition~\ref{prop:rule-exact}\,(vi)) the endpoint is still admissible by
(i) and only this bound is unavailable; at $\kappa^{+}=n-1$ the excess
$\tilde k^{+}-\kappa^{+}$ is at most one count.
\end{enumerate}
\end{proposition}

\begin{proof}
\emph{(i).} $\mathsf{C}^{+}(\tilde k^{+})$ gives
$\Pr[\Ncount{}>\tilde k^{+}]\le\tau^{+}(\tilde k^{+})\le\varepsilon^{+}$, or
in log space $\Pr[\Ncount{}>\tilde k^{+}]\le e^{U(\tilde k^{+})}
\le e^{\lambda^{-}}\le\varepsilon^{+}$. So $\tilde k^{+}$ belongs to the set
in \eqref{eq:kappaplus}, and $\kappa^{+}\le\tilde k^{+}$ by the minimality
that defines $\kappa^{+}$ (Proposition~\ref{prop:rule-exact}\,(iii)). The
lower side is the same argument for $\Ncount{}'$: a certified index
$\tilde k'$ for the upper tail of $\Ncount{}'$ at level $\varepsilon^{-}$
gives $\Pr[\Ncount{}<n-\tilde k']=\Pr[\Ncount{}'>\tilde k']\le\varepsilon^{-}$
by Lemma~\ref{lem:tail-majorant}\,(v), so $\tilde k^{-}\coloneqq n-\tilde k'$
lies in the set of \eqref{eq:kappaminus} and $\tilde k^{-}\le\kappa^{-}$ by
maximality. Both values lie in $\{0,\dots,n\}$; with the conventions at a
zero budget and at $p\in\{0,1\}$ imposed by definition, $R_{\mathrm{encl}}$
satisfies Definition~\ref{def:endpoint-rule}\,(a)--(c). For the containment:
by \eqref{eq:guard} the \tref{guarded endpoints} are monotone in the rule's
endpoints, so $\tilde k^{-}\le\kappa^{-}$ and $\tilde k^{+}\ge\kappa^{+}$
give $k^{-}_{\mathrm{encl}}(\hat c)\le k^{-}_{\mathrm{ex}}(\hat c)$ and
$k^{+}_{\mathrm{encl}}(\hat c)\ge k^{+}_{\mathrm{ex}}(\hat c)$ at every
constrained outcome, hence the same order for the bounds \eqref{eq:boxbounds}
and the inclusion of the sets \eqref{eq:accset}.

\emph{(ii).} If $\tilde k^{+}=0$ then $\kappa^{+}\le0$ by (i), so
$\kappa^{+}=0$. Otherwise $\mathsf{D}^{+}(\tilde k^{+}-1)$ gives
$\Pr[\Ncount{}>\tilde k^{+}-1]\ge\tau^{-}(\tilde k^{+}-1)>\varepsilon^{+}$
(in log space
$\Pr[\Ncount{}>\tilde k^{+}-1]\ge e^{\Lambda(\tilde k^{+}-1)}
>e^{\lambda^{+}}\ge\varepsilon^{+}$), so $\tilde k^{+}-1$ is not in the set
of \eqref{eq:kappaplus}. That set is an up-set of $\{0,\dots,n\}$, because
$k\mapsto\Pr[\Ncount{}>k]$ is non-increasing (Lemma~\ref{lem:guard}\,(ii));
hence its minimum $\kappa^{+}$ exceeds $\tilde k^{+}-1$, i.e.\
$\kappa^{+}\ge\tilde k^{+}$, and with (i) equality. The lower side is the
mirror.

\emph{(iii).} Since $\tilde k^{+}>\kappa^{+}\ge0$, the index $\tilde k^{+}-1$
exists and is not certified admissible, so
$\tau^{+}(\tilde k^{+}-1)>\varepsilon^{+}$ (in log space
$U(\tilde k^{+}-1)>\lambda^{-}$), and by the definition of $w$,
\begin{equation*}
  \ln\Pr[\Ncount{}>\tilde k^{+}-1]\ =\ \ln\tau^{+}(\tilde k^{+}-1)-w
  \ >\ \ln\varepsilon^{+}-w
\end{equation*}
(in log space $\ln\Pr[\Ncount{}>\tilde k^{+}-1]=U(\tilde k^{+}-1)
-\big(U(\tilde k^{+}-1)-\ln\Pr[\Ncount{}>\tilde k^{+}-1]\big)
>\lambda^{-}-\big(w-(\ln\varepsilon^{+}-\lambda^{-})\big)
=\ln\varepsilon^{+}-w$). On the other hand $\tilde k^{+}-1\ge\kappa^{+}$,
and for every $l$ with $\kappa^{+}+1\le l\le\tilde k^{+}-1$ one has
$r_l\le r_{\kappa^{+}+1}<1$ by Lemma~\ref{lem:tail-majorant}\,(i), so
\eqref{eq:tail-decay} gives $\Pr[\Ncount{}>l]\le r_l\Pr[\Ncount{}>l-1]
\le r_{\kappa^{+}+1}\Pr[\Ncount{}>l-1]$; iterating from $l=\kappa^{+}+1$ to
$l=\tilde k^{+}-1$,
\begin{equation*}
  \Pr[\Ncount{}>\tilde k^{+}-1]\ \le\
  r_{\kappa^{+}+1}^{\,\tilde k^{+}-1-\kappa^{+}}\,\Pr[\Ncount{}>\kappa^{+}]
  \ \le\ r_{\kappa^{+}+1}^{\,\tilde k^{+}-1-\kappa^{+}}\,\varepsilon^{+},
\end{equation*}
the last step by \eqref{eq:exact-admissible}. Taking logarithms and
combining with the previous display,
$\ln\varepsilon^{+}-w<\ln\varepsilon^{+}+(\tilde k^{+}-1-\kappa^{+})
\ln r_{\kappa^{+}+1}$, i.e.\
$(\tilde k^{+}-1-\kappa^{+})(-\ln r_{\kappa^{+}+1})<w$, which is
\eqref{eq:bracket-conservatism} since $-\ln r_{\kappa^{+}+1}$ is positive and
finite: $r_{\kappa^{+}+1}<1$ by hypothesis, and $r_{\kappa^{+}+1}>0$ by
Lemma~\ref{lem:tail-majorant}\,(i) because $\kappa^{+}+1\le n-1$. The mirror
statement is the same argument for the upper tail of $\Ncount{}'$ at level
$\varepsilon^{-}$, translated by $k\mapsto n-k$: the exact lower quantile of
$\Ncount{}$ is $n$ minus the exact upper quantile of $\Ncount{}'$, because
$\Pr[\Ncount{}<k]=\Pr[\Ncount{}'>n-k]$ turns the maximum in
\eqref{eq:kappaminus} into the minimum in \eqref{eq:kappaplus} for
$\Ncount{}'$.

\emph{The closing sentences of (iii).} By \eqref{eq:guard} the guarded upper
endpoint is $\max\{\tilde k^{+},\lceil np\rceil\}$ for this evaluation and
$\max\{\kappa^{+},\lceil np\rceil\}$ for the exact rule, and by
\eqref{eq:boxbounds} the box bound is the guarded endpoint divided by $n$.
Where the bracket is decided, (ii) gives $\tilde k^{+}=\kappa^{+}$, so the two
guarded endpoints and the two bounds coincide. Otherwise, the map
$x\mapsto\max\{x,\lceil np\rceil\}$ is non-decreasing and changes by at most
$x'-x$ between $x\le x'$, so with $\kappa^{+}\le\tilde k^{+}$ from (i) the
guarded endpoint of this evaluation exceeds the exact rule's by at least $0$
and at most $\tilde k^{+}-\kappa^{+}$, and dividing by $n$ gives the bound on
$q^{+}$; \eqref{eq:bracket-conservatism} bounds $\tilde k^{+}-\kappa^{+}$
under the hypotheses of (iii). At $\kappa^{+}=n-1$, $\tilde k^{+}\le n$ gives
$\tilde k^{+}-\kappa^{+}\le1$. The lower side is the mirror, with
$x\mapsto\min\{x,\lfloor np\rfloor\}$ and $\tilde k^{-}\le\kappa^{-}$.
\end{proof}
